%% Exam randomization
\RequirePackage{tikz,pgf,siunitx,etoolbox,xparse,xintfrac,xintexpr}
\usetikzlibrary{math}

\ExplSyntaxOn%
% %Claude.ai was used to help generate the \randNums macro
%%%%%%%%%%%%%%%%%%%%%%%%%%%%%%%%%%%%%%%%%%%%%%%%%%%%%%%%%%%%%%%%%%%%%%%%%%%%%%%%
\int_new:N    \l__dp_lower_int   % decimal places in lower bound string
\int_new:N    \l__dp_upper_int   % decimal places in upper bound string
\seq_new:N    \l__dp_seq         % scratch for splitting on "."
\tl_new:N     \l__dp_frac_tl     % scratch for the fractional part

%% count num decimal places on bounds
\cs_new_protected:Npn \__randints_count_dp:NN #1 #2
{
  \tl_set:Nx \l__dp_frac_tl { \tl_use:N #2 }
  % Split on ".": one item = integer, two items = has decimal part
  \seq_set_split:NnV \l__dp_seq { . } \l__dp_frac_tl
  \int_compare:nNnTF { \seq_count:N \l__dp_seq } = { 1 }
  { \int_set:Nn #1 { 0 } }
  {
    \tl_set:Nx \l__dp_frac_tl { \seq_item:Nn \l__dp_seq { 2 } }
    \int_set:Nn #1 { \tl_count:N \l__dp_frac_tl }
  }
}

%%%%%%%%%%%%%%%%%%%%%%%%%%%%%%%%%%%%%%%%%%%%%%%%%%%%%%%%%%%%%%%%%%%%%%%%%%%%%%%%
\bool_new:N   \l__pick_found_bool      % track if rand num in skip
\bool_new:N   \l__randNums_uniq_bool   % [uniq] or star
\int_new:N    \l__randNums_min_int     % sorted min
\int_new:N    \l__randNums_max_int     % sorted max
\int_new:N    \l__pool_size_int        % stores size of pool of ints
\int_new:N    \l__sort_idx_int         % index for sorting values
\seq_new:N    \l__sort_seq             % sequence to hold values for sorting
\tl_new:N     \l__my_pick_tl           % scratch for randomly chosen value

\keys_define:nn { randNums }
  {
    nz            .bool_set:N       = \l__randNums_kv_nonzero_bool,
    nz            .initial:n        = false,
    nz            .default:n        = true,
    %
    sort          .tl_set:N         = \l__randNums_kv_sort_tl,
    sort          .initial:n        = \q_no_value,
    sort          .default:n        = >,
    %
    uniq          .bool_set:N       = \l__randNums_kv_uniq_bool,
    uniq          .initial:n        = false,
    uniq          .default:n        = true,
    %
    skip          .clist_set:N      = \l__randNums_kv_skip_clist,
    %
    min           .code:n           = {
                    \fp_set:Nn  \l__randNums_kv_min_fp { #1 }
                    \tl_set:Nn  \l__randNums_kv_min_tl { #1 } },
    min           .initial:n        = 1,
    min           .value_required:n = true,
    %
    max           .code:n           = {
                    \fp_set:Nn  \l__randNums_kv_max_fp { #1 }
                    \tl_set:Nn  \l__randNums_kv_max_tl { #1 } },
    max           .initial:n        = 5,
    max           .value_required:n = true,
    %
    mult          .fp_set:N         = \l__randNums_kv_mult_fp,
    mult          .initial:n        = 1,
    mult          .value_required:n = true,
    %
    numDec        .fp_set:N         = \l__randNums_kv_numDec_fp,
    numDec        .initial:n        = nan,
    numDec        .value_required:n = true,
  }

%%%%%%%%%%%%%%%%%%%%%%%%%%%%%%%%%%%%%%%%%%%%%%%%%%%%%%%%%%%%%%%%%%%%%%%%%%%%%%%%
%%%%%%%%%%%%%%%%%%%%%%%%%%%%%%%% Error messages %%%%%%%%%%%%%%%%%%%%%%%%%%%%%%%%
%%%%%%%%%%%%%%%%%%%%%%%%%%%%%%%%%%%%%%%%%%%%%%%%%%%%%%%%%%%%%%%%%%%%%%%%%%%%%%%%
\msg_new:nnn { randNums } { too-many }
  { \randNums*:~only~#1~value(s)~available~in~[#3]~but~#2~requested. }

% \randNums * [keys] {csname1,csname2,...}
% #1 - starred variant produces unique values
% #2 - optional args: (default)
%          nz = (false) generate nonzero values from [-max,-min]U[min,max]
%        sort = (<) sort random values so that csname1<csname2<...
%        uniq = (false) produce unique values
%        skip = () clist of values to skip
%         min = (1) minimum value of range
%         max = (5) maximum value of range
%        mult = (1) value to multiply by after drawing (e.g. random multiples of 5)
%      numDec = (0) sets number of decimal places in sample space (e.g. 1, 1.1, 1.2, ..., 5)
% #3 - clist of names; defines \name globally to the drawn value
\NewDocumentCommand \randNums { s O{} m }{
  \group_begin: %grouping used to keep key changes local
  \keys_set:nn { randNums }{ #2 }

  % Star is equivalent to uniq
  \bool_if:nT {#1} { \bool_set_true:N \l__randNums_kv_uniq_bool }

  % Set uniq_bool based on keyval pair
  \bool_set_false:N \l__randNums_uniq_bool
  \bool_if:NT \l__randNums_kv_uniq_bool { \bool_set_true:N \l__randNums_uniq_bool }

  % % Detect decimal places in each bound's string form
  \__randints_count_dp:NN \l__dp_lower_int { \l__randNums_kv_min_tl }
  \__randints_count_dp:NN \l__dp_upper_int { \l__randNums_kv_max_tl }

  % Scale factor = 10^(max decimal places); 1 when both bounds are integers
  \fp_set:Nn \l__dp_scale_fp
  { \fp_eval:n { 10 ^ \int_max:nn { \l__dp_lower_int } { \l__dp_upper_int } } }

  %manually override scaling
  \fp_if_nan:nF { \l__randNums_kv_numDec_fp }{
    \fp_set:Nn \l__dp_scale_fp { \fp_eval:n { 10 ^ \l__randNums_kv_numDec_fp }}
    }

  % set the upper and lower bounds for the random integers
  \int_set:Nn \l__randNums_min_int {
    \fp_eval:n{ round(min(
      \l__randNums_kv_min_fp * \l__dp_scale_fp ,
      \l__randNums_kv_max_fp * \l__dp_scale_fp))
    }
  }
  \int_set:Nn \l__randNums_max_int {
    \fp_eval:n{ round(max(
      \l__randNums_kv_min_fp * \l__dp_scale_fp,
      \l__randNums_kv_max_fp * \l__dp_scale_fp))
    }
  }

  %build a sequence of "used values"
  \seq_clear_new:N \l__used_seq
  \clist_map_inline:Nn \l__randNums_kv_skip_clist {
    \bool_if:NTF { \l__randNums_kv_nonzero_bool }{
      % nz pool: ±[lower..upper]*multiplier, so check |value| is in range
      \fp_compare:nNnTF { abs(##1) } < { \l__randNums_min_int * \l__randNums_kv_mult_fp }{
        % out of range low - skip
        }{
          \fp_compare:nNnTF { abs(##1) } > { \l__randNums_max_int * \l__randNums_kv_mult_fp }{
            % out of range high - skip
          }{
            \seq_put_right:Nx \l__used_seq { \fp_eval:n { ##1 } }
          }
      }
    }{
      % normal pool: check lower*multiplier <= value <= upper*multiplier
      \fp_compare:nNnTF { ##1 } < {
        \l__randNums_min_int
        / \l__dp_scale_fp
        * \l__randNums_kv_mult_fp
        }
        { } % out of range - skip
        {
          \fp_compare:nNnTF { ##1 } > {
            \l__randNums_max_int
            / \l__dp_scale_fp
            * \l__randNums_kv_mult_fp
          }{
            % out of range - skip
          }{
            \seq_put_right:Nx \l__used_seq { \fp_eval:n { ##1 } }
          }
        }
    }
  }

  % overflow check (only meaningful without replacement)
  \bool_if:NTF { \l__randNums_uniq_bool }{
    \bool_if:NTF { \l__randNums_kv_nonzero_bool }
      { \int_set:Nn \l__pool_size_int { 2 * ( \l__randNums_max_int - \l__randNums_min_int + 1 ) } }
      { \int_set:Nn \l__pool_size_int {       \l__randNums_max_int - \l__randNums_min_int + 1   } }
    % Subtract the number of values being skipped
    \int_sub:Nn \l__pool_size_int { \seq_count:N \l__used_seq }
    % Compare against the number of values requested
    \int_compare:nNnT { \clist_count:n { #3 } } > { \l__pool_size_int }{
      \msg_fatal:nnxxx { randNums } { too-many }
        { \int_use:N \l__pool_size_int }
        { \clist_count:n { #3 } }
        { \bool_if:NTF { \l__randNums_kv_nonzero_bool } {nz~}
          {} \int_eval:n { \l__randNums_min_int }
          .. \int_eval:n { \l__randNums_max_int }}%
    }
  }{
  }

  %generate random integers
  \clist_map_inline:nn {#3}{
    \bool_set_false:N \l__pick_found_bool
    \bool_until_do:Nn \l__pick_found_bool{
      \bool_if:NTF { \l__randNums_kv_nonzero_bool }{
        \tl_set:Nx \l__my_pick_tl{
          \fp_eval:n{
            \int_rand:nn {1}{2} == 1
              ? \int_rand:nn
                {\l__randNums_min_int}
                {\l__randNums_max_int}
                / \l__dp_scale_fp
                * \l__randNums_kv_mult_fp
              : -\int_rand:nn
                {\l__randNums_min_int}
                {\l__randNums_max_int}
                / \l__dp_scale_fp
                * \l__randNums_kv_mult_fp
            }
          }
        }{
          \tl_set:Nx \l__my_pick_tl
            { \fp_eval:n { \int_rand:nn
                {\l__randNums_min_int}
                {\l__randNums_max_int}
                / \l__dp_scale_fp
                * \l__randNums_kv_mult_fp
              } }
        }
      % Accept the pick only if it's not in the exclusion list
      \seq_if_in:NVF \l__used_seq \l__my_pick_tl { \bool_set_true:N \l__pick_found_bool }
    }

    %Add value generated to used sequence if done w/o replacement
    \bool_if:NTF { \l__randNums_uniq_bool }{ \seq_put_right:NV \l__used_seq \l__my_pick_tl }{}
    %Globally store value generated
    \cs_gset:cpx {##1}{ \tl_use:N \l__my_pick_tl }
  }

  %sort values
  \quark_if_no_value:NF \l__randNums_kv_sort_tl {
    \seq_clear:N \l__sort_seq
    \clist_map_inline:nn {#3}
      { \seq_put_right:Nx \l__sort_seq { \use:c { ##1 } } }   % \use:c dereferences
    \seq_sort:Nn \l__sort_seq {
      \str_case:en { \l__randNums_kv_sort_tl }
      {
        { < } { \fp_compare:nNnTF { ##2 } < { ##1 }
                  { \sort_return_swapped: } { \sort_return_same: } }
        { > } { \fp_compare:nNnTF { ##2 } > { ##1 }
                  { \sort_return_swapped: } { \sort_return_same: } }
      }
    }
    \int_zero:N \l__sort_idx_int
    \clist_map_inline:nn {#3} {
      \int_incr:N \l__sort_idx_int
      \cs_gset:cpx {##1} { \seq_item:Nn \l__sort_seq { \l__sort_idx_int } }
    }
  }
  \group_end:
}
\ExplSyntaxOff%
\RequirePackage{siunitx}
\ExplSyntaxOn%
\cs_if_exist:NF \fp_format:nn { \RequirePackage{l3str-format} } % needed for TL2025
% \debug_on:n { check-declarations }
%%%%%%%%%%%%%%%%%%%%%%%%%%%%%%%%%%%%%%%%%%%%%%%%%%%%%%%%%%%%%%%%%%%%%%%%%%%%%%%%
%\num from siunitx for explicit printing of +
\cs_new_protected:Npn \__printcoeff_num:nn #1 #2
  { \num[#1]{#2} }

\bool_new:N   \l__printFuncs_printSign_bool
\fp_new:N     \l__tmpa_fp
\int_new:N    \l__tmpa_int

\keys_define:nn { printNum }
  {
    format      .tl_set:N         =  \l__printcoeff_fmt_tl ,
    format      .value_required:n = true,
    %
    fmt         .meta:n           = { format = #1 }, %alias for format
    fmt         .value_required:n = true,
    %
    noplus      .bool_set:N       =  \l__printcoeff_noplus_bool ,
    noplus      .initial:n        =  false ,
    noplus      .default:n        =  true ,  %set when called without setting
    %
    numDec      .int_set:N        = \l_printcoeff_numDec_int,
    numDec      .initial:n        = 4,
    numDec      .value_required:n = true,
    %
    signs-only  .bool_set:N       =  \l__signonly_units_bool ,
    signs-only  .initial:n        =  true ,
    signs-only  .default:n        =  false , %set when called without setting
    %
    one         .meta:n           =  { signs-only = #1 }, %alias for `signs-only'
    one         .initial:n        =  true ,
    one         .default:n        =  false , %set when called without setting
    %
    zero        .bool_set:N       =  \l__print_zero_bool ,
    zero        .initial:n        =  false ,
    zero        .default:n        =  true ,  %set when called without setting
  }

% \printNum * [keys] {<coeff>}
% #1 - starred variant suppress leading '+' (equivalent to noplus key)
% #2 - optional args: (default)
%       format=<fmt>: ()      fp_format string (default: empty)
%             noplus: (false) toggle suppressing leading '+' without the star
%             numdec: (4)     number of decimal places to round to
%                one: (true)  toggle printing one (alias for signs-only)
%               zero: (false) toggle printing zero
\NewDocumentCommand \printNum { s O{} m }
  { \print_coeff:nnn {#1} {#2} {#3} }

\cs_new_protected:Npn \print_coeff:nnn #1 #2 #3
  {
    \group_begin:
      \keys_set:nn { printNum } { #2 }

      % Star is equivalent to noplus
      \bool_if:nT {#1} { \bool_set_true:N \l__printcoeff_noplus_bool }

      % Derive print_sign from noplus
      \bool_set_true:N \l__printFuncs_printSign_bool
      \bool_if:NT \l__printcoeff_noplus_bool
        { \bool_set_false:N \l__printFuncs_printSign_bool }

      \fp_set:Nn \l__tmpa_fp { \fp_eval:n { #3 } } % numeric form for comparisons

      \bool_if:nTF {
        \fp_compare_p:n {abs(\l__tmpa_fp)==1} && \l__signonly_units_bool
      }{
        \fp_compare:nTF {\l__tmpa_fp == -1}
          {\ensuremath{-}}  %-1 w/ signonly_units
          {
           \bool_if:NTF \l__printFuncs_printSign_bool
           {\ensuremath{+}} %1 w/  sign
           {}               %1 w/o sign
         }
      }{

        \bool_if:nT {% print zero or non-zero number: generic case
          \l__print_zero_bool || \fp_compare_p:n {\l__tmpa_fp != 0}
        }
        {
          \group_begin:
          \ensuremath{
            \exp_args:Ne \__printcoeff_num:nn
            {
              group-separator={,},
              group-digits=integer,
              \bool_if:NT \l__printFuncs_printSign_bool { print-implicit-plus=true }
            }
            {
              \exp_args:NnV \fp_format:nn
              { round(\l__tmpa_fp,\l_printcoeff_numDec_int) } \l__printcoeff_fmt_tl
            }
          }
          \group_end:
        }
      }
    \group_end:
    }

%%%%%%%%%%%%%%%%%%%%%%%%%%%%%%%%%%%%%%%%%%%%%%%%%%%%%%%%%%%%%%%%%%%%%%%%%%%%%%%%
\bool_new:N   \l__printpoly_coeffOne_bool
\bool_new:N   \l__printpoly_expozero_bool
\bool_new:N   \l__printpoly_firstTerm_bool
\fp_new:N     \l__printpoly_coeff_fp
\fp_new:N     \l__printpoly_expo_fp
\int_new:N    \l__printpoly_expoidx_int
\seq_new:N    \l__printpoly_expo_seq
\tl_new:N     \l__printpoly_coeff_tl

\keys_define:nn { printPoly }
  {
    deriv       .int_set:N        = \l__printpoly_deriv_int,
    deriv       .initial:n        = 0,
    deriv       .default:n        = 1,
    %
    format      .tl_set:N         = \l__printpoly_fmt_tl,
    format      .value_required:n = true,
    %
    fmt         .meta:n           = { format = #1 }, %alias for format
    fmt         .value_required:n = true,
    %
    sort        .tl_set:N         = \l__printPoly_sort_tl,
    sort        .initial:n        = >,
    sort        .value_required:n = true,
    %
    var         .tl_set:N         = \l__printpoly_var_tl,
    var         .initial:n        = x,
    var         .value_required:n = true,
  }

% \printPoly * [keys] {<coeff>}{<expo>}
% #1 - starred variant enable leading '+'
% #2 - optional args: (default)
%              deriv: (0) number of derivative using power rule
%       format=<fmt>: ()  fp_format string
%               sort: (>) sort auto-generated expo args
%                var: (x) variable used in expression
% #3 - clist of coefficients
% #4 - (optional) clist of exponents
\NewDocumentCommand \printPoly { s O{} m m }
  { \print_poly:nnnn {#1} {#2} {#3} {#4} }

\cs_new_protected:Npn \print_poly:nnnn #1 #2 #3 #4
{
  \group_begin:
    \keys_set:nn { printPoly } { #2 }

    \bool_set_true:N \l__printpoly_firstTerm_bool
    \bool_if:nT {#1} { \bool_set_false:N \l__printpoly_firstTerm_bool }

    % iterate over #3, store #4 as a sequence
    \seq_set_from_clist:Nn \l__printpoly_expo_seq { #4 }
    %build sequence if #4 is empty
    \tl_if_blank:nT { #4 } {
      \seq_clear:N \l__printpoly_expo_seq
      \int_step_inline:nnnn { 0 } { 1 }  { \clist_count:n {#3} - 1 } { \seq_put_right:Nn \l__printpoly_expo_seq {##1} }
      %logic to reverse order
      \str_case:en { \l__printPoly_sort_tl }
        {
          { < } { }
          { > } { \seq_reverse:N \l__printpoly_expo_seq }
        }
      }

    %TODO: error message if #3 and #4 are different lengths

    \clist_map_inline:nn { #3 }{
      \int_incr:N \l__printpoly_expoidx_int
      % get current coeff and expo
      \fp_set:Nn \l__printpoly_coeff_fp { ##1 }
      \fp_set:Nn \l__printpoly_expo_fp {
          \seq_item:Nn \l__printpoly_expo_seq { \l__printpoly_expoidx_int } }

      %derivative power rule
      \int_step_inline:nn { \l__printpoly_deriv_int }{
        \int_set:Nn \l__tmpa_int { ####1 }
        \fp_set:Nn \l__tmpa_fp { \l__printpoly_expo_fp+1-\l__tmpa_int }
        \fp_set:Nn \l__printpoly_coeff_fp { \l__printpoly_coeff_fp*(\l__tmpa_fp) }
      }
      % VERY rudimentary integral power rule
      \int_step_inline:nn { -\l__printpoly_deriv_int }{
        \int_set:Nn \l__tmpa_int { ####1 }
        \fp_set:Nn \l__tmpa_fp { \l__printpoly_expo_fp+2-\l__tmpa_int }
        \fp_compare:nNnF { \l__tmpa_fp } = { 0 }
          { \fp_set:Nn \l__printpoly_coeff_fp { \l__printpoly_coeff_fp/(\l__tmpa_fp) }}
      }
      %modify exponent
      \fp_set:Nn \l__printpoly_expo_fp {\l__printpoly_expo_fp - \l__printpoly_deriv_int}

      \fp_compare:nNnF { \l__printpoly_coeff_fp } = {0}{
        \tl_set:Nn \l__printpoly_coeff_tl { \l__printpoly_coeff_fp }

        \bool_set:Nn \l__printpoly_expozero_bool {\fp_compare_p:nNn { \l__printpoly_expo_fp } = {0}}
        \bool_set:Nn \l__printpoly_coeffOne_bool {\fp_compare_p:nNn { abs(\l__printpoly_coeff_fp) } = {1}}

        \bool_case:n
        {
          { \bool_lazy_and_p:nn  % first term, cx
              {  \l__printpoly_firstTerm_bool }
              { !\l__printpoly_expozero_bool }}
            { \printNum*{ \l__printpoly_coeff_tl}
              \bool_set_false:N \l__printpoly_firstTerm_bool }
          { \bool_lazy_and_p:nn  % first term, c
              {  \l__printpoly_firstTerm_bool }
              {  \l__printpoly_expozero_bool }}
            { \printNum*[one]{ \l__printpoly_coeff_tl}
              \bool_set_false:N \l__printpoly_firstTerm_bool }
          { \bool_lazy_all_p:n { % not first term, \pm 1
              { !\l__printpoly_firstTerm_bool }
              {  \l__printpoly_expozero_bool }
              {  \l__printpoly_coeffOne_bool }}}
            { \printNum[one]{ \l__printpoly_coeff_tl} }
          { \bool_lazy_and_p:nn  % not first term, \pm c or \pm cx
              { !\l__printpoly_firstTerm_bool }
              { !\l__printpoly_expozero_bool ||
                !\l__printpoly_coeffOne_bool }}
            { \printNum{ \l__printpoly_coeff_tl} }
        }

        \fp_compare:nNnF { \l__printpoly_expo_fp } = {0}{
          %Print if expo nonzero
          \ensuremath{\l__printpoly_var_tl}
          %Print expo if not one
          \fp_compare:nNnF { \l__printpoly_expo_fp } = {1}{
            \ensuremath{^{ \printNum[noplus,one]{\l__printpoly_expo_fp }}}
          }
        }
      }
  }
  \group_end:
}

\ExplSyntaxOff%

%TODO: remove these lengths?
\newcommand{\ansBoxHeight}{\ifprintanswers8mm\else11mm\fi}
\newcommand{\ansBoxWidth}{35mm}

\ExplSyntaxOn
% \debug_on:n { check-declarations }

\keys_define:nn { ansBox }
  {
    height      .dim_set:N   = \l__ansBox_height_dim,
    height      .initial:n   = \ifprintanswers8mm\else11mm\fi,

    width       .dim_set:N   = \l__ansBox_width_dim,
    width       .initial:n   = 35mm,
  }
% creates answer box of (opt) height, (opt) width, with mandatory contents
% \ansBox [keys] {contents}
% #1 - optional args: (default)
%     height = (\ifprintanswers8mm\else11mm\fi) sets height of ansBox
%      width = (35mm) sets width of ansBox
\NewDocumentCommand{\ansBox}{ O{} m }{%
  \keys_set:nn { ansBox }{ #1 }
  \hspace*{\stretch{1}}\fbox{\begin{minipage}[t][\l__ansBox_height_dim][c]{\l__ansBox_width_dim}#2\end{minipage}}\newline%
}

%%%%%%%%%%%%%%%%%%%%%%%%%%%%%%%%%%%%%%%%%%%%%%%%%%%%%%%%%%%%%%%%%%%%%%%%%%%%%%%%
% Manually define a variable (control sequence)
\NewDocumentCommand \defVar { m m }{
  \cs_gset:cpx { #1 }{ \fp_eval:n { #2 }}
}
%%%%%%%%%%%%%%%%%%%%%%%%%%%%%%%%%%%%%%%%%%%%%%%%%%%%%%%%%%%%%%%%%%%%%%%%%%%%%%%%
%TODO \NewDocumentCommand \sortVars { m m }{
% sorts a list of values and stores in a list of csnames
% #1 - clist of csnames
% #2 - clist of values to sort
%}
%%%%%%%%%%%%%%%%%%%%%%%%%%%%%%%%%%%%%%%%%%%%%%%%%%%%%%%%%%%%%%%%%%%%%%%%%%%%%%%%
\NewDocumentCommand \defMinMax { m m m m}{
  \cs_gset:cpx { #1 }{ \fp_min:nn { #3 }{ #4 }}
  \cs_gset:cpx { #2 }{ \fp_max:nn { #3 }{ #4 }}
}

%%%%%%%%%%%%%%%%%%%%%%%%%%%%%%%%%%%%%%%%%%%%%%%%%%%%%%%%%%%%%%%%%%%%%%%%%%%%%%%%
\NewDocumentCommand \randItem { m m }{
  \cs_gset:cpx { #1 }{ \clist_rand_item:n { #2 }}
}

%%%%%%%%%%%%%%%%%%%%%%%%%%%%%%%%%%%%%%%%%%%%%%%%%%%%%%%%%%%%%%%%%%%%%%%%%%%%%%%%
%% \fpIfGreater{a}{b}{true}{false}
\NewDocumentCommand \fpIfGreater { m m m m }
  {
    \fp_compare:nNnTF {#1} > {#2} {#3} {#4}
  }

%%%%%%%%%%%%%%%%%%%%%%%%%%%%%%%%%%%%%%%%%%%%%%%%%%%%%%%%%%%%%%%%%%%%%%%%%%%%%%%%
% Print negative numbers inside parenthesis
\NewDocumentCommand \printWithParens { m }{
  \fp_compare:nNnTF {#1} < {0}{
    \ensuremath{(#1)}
  }{
    \ensuremath{#1}
  }
}
%%%%%%%%%%%%%%%%%%%%%%%%%%%%%%%%%%%%%%%%%%%%%%%%%%%%%%%%%%%%%%%%%%%%%%%%%%%%%%%%
\bool_new:N \l_isprime_bool
\int_new:N \l_trial_int

\NewDocumentCommand \IfPrime { m m m } {
  % #1 = number to test, #2 = true code, #3 = false code
  \bool_set_true:N \l_isprime_bool

  \int_compare:nNnT { #1 } < { 2 } { \bool_set_false:N \l_isprime_bool }

  \int_if_even:nT { #1 } {
    \int_compare:nNnT { #1 } > { 2 } { \bool_set_false:N \l_isprime_bool }
  }

  % Trial division by odd numbers from 3 up to floor(sqrt(#1))
  \int_set:Nn \l_trial_int { 3 }
  \bool_while_do:nn
    {
      \l_isprime_bool
      &&
      \int_compare_p:n { \l_trial_int * \l_trial_int <= #1 }
    }
    {
      \int_compare:nNnT { \int_mod:nn{#1}{\l_trial_int} } = { 0 }
        { \bool_set_false:N \l_isprime_bool }
      \int_add:Nn \l_trial_int { 2 }
    }

  \bool_if:NTF \l_isprime_bool { #2 } { #3 }
}

\NewDocumentCommand \randPrime { m m m } {
  % #1 control sequence, #2 = lower bound, #3 = upper bound
  % Result stored in \l_candidate_int (and printed)
  \bool_set_false:N \l_isprime_bool
  \bool_until_do:Nn \l_isprime_bool {
    \tl_set:Nx \l_tmpa_tl { \int_rand:nn{ #2 }{ #3 } }
    \IfPrime{ \l_tmpa_tl }
      { \bool_set_true:N \l_isprime_bool }
      {}
  }
  \cs_gset:cpx { #1 }{ \tl_use:N \l_tmpa_tl }
}
%%%%%%%%%%%%%%%%%%%%%%%%%%%%%%%%%%%%%%%%%%%%%%%%%%%%%%%%%%%%%%%%%%%%%%%%%%%%%%%%
%TODO: Add optional argument for separators when printing sequence
\NewDocumentCommand{\ShuffleList}{m}{
  % Build a sequence from the input list
  \seq_clear:N \l_shuffle_seq
  \clist_map_inline:nn {#1} {
    \seq_put_right:Nn \l_shuffle_seq {##1}
  }
  % Shuffle sequence
  \seq_gshuffle:N \l_shuffle_seq

  \seq_use:Nnnn \l_shuffle_seq {}{}{}
}
%%%%%%%%%%%%%%%%%%%%%%%%%%%%%%%%%%%%%%%%%%%%%%%%%%%%%%%%%%%%%%%%%%%%%%%%%%%%%%%%
% Creates and sorts sequences
\cs_new_protected:Npn \__sortintseq_add:Nn #1 #2{
    \seq_gclear:N \g__sortintseq_tmp_seq
    \clist_map_inline:nn { #2 }
      { \seq_gput_right:No \g__sortintseq_tmp_seq { ##1 } }
    \seq_gsort:Nn \g__sortintseq_tmp_seq {
      \int_compare:nNnTF { ##1 } > { ##2 }
        { \sort_return_swapped: }
        { \sort_return_same: }
    }
    \cs_gset_eq:NN #1 \g__sortintseq_tmp_seq  % assign result to user's variable
  }

\NewDocumentCommand \SortIntSeq { mm }{
    \__sortintseq_add:Nn #1 { #2 }
}
%Print a sequence
\NewDocumentCommand \PrintSeq { m }{
    \seq_use:Nn #1 { ,~ }
}

%%%%%%%%%%%%%%%%%%%%%%%%%%%%%%%%%%%%%%%%%%%%%%%%%%%%%%%%%%%%%%%%%%%%%%%%%%%%%%%%
\NewDocumentCommand \printIfTrue { m m }{
  \ifnum #1 #2 \fi
  }
%%%%%%%%%%%%%%%%%%%%%%%%%%%%%%%%%%%%%%%%%%%%%%%%%%%%%%%%%%%%%%%%%%%%%%%%%%%%%%%%
\NewDocumentCommand \printIfFalse { m m }{
  \ifnum #1 \else #2 \fi
  }
%%%%%%%%%%%%%%%%%%%%%%%%%%%%%%%%%%%%%%%%%%%%%%%%%%%%%%%%%%%%%%%%%%%%%%%%%%%%%%%%
\NewDocumentCommand \defGCD { m m m }{
  \cs_gset:cpx { #1 }{ \exam_gcd:nn #2 #3 }
}

\cs_new:Npn \exam_gcd:nn #1 #2
  {
    \int_eval:n {
        \int_compare:nTF { #2 = 0 }
          { #1 }
          { \exam_gcd:nn { #2 } { \int_mod:nn { \int_abs:n{#1} } { \int_abs:n{#2} } } }
      }
  }

\NewDocumentCommand \defLCM { m m m }{
  \cs_gset:cpx { #1 }{ \int_lcm:nn #2 #3 }
}

\cs_new:Npn \int_lcm:nn #1 #2
  {
    \int_eval:n { ( \int_abs:n{#1} / \exam_gcd:nn { #1} {#2} ) * \int_abs:n{#2} }
  }
%%%%%%%%%%%%%%%%%%%%%%%%%%%%%%%%%%%%%%%%%%%%%%%%%%%%%%%%%%%%%%%%%%%%%%%%%%%%%%%%
% Translates 1,...,26 to corresponding capital letter
\NewDocumentCommand \numToLetter { m }{
  \int_to_Alph:n { #1 }
}

\NewDocumentCommand \numToletter { m }{
  \int_to_alph:n { #1 }
}
%%%%%%%%%%%%%%%%%%%%%%%%%%%%%%%%%%%%%%%%%%%%%%%%%%%%%%%%%%%%%%%%%%%%%%%%%%%%%%%%
% Translates letters to numbers
\NewDocumentCommand \letterToNum { m m }{
  \cs_gset:cpx { # 1}{ \int_from_alph:n{#2} }
}
%%%%%%%%%%%%%%%%%%%%%%%%%%%%%%%%%%%%%%%%%%%%%%%%%%%%%%%%%%%%%%%%%%%%%%%%%%%%%%%%
\NewDocumentCommand{\signOnly}{m}{%
  \printNum[signs-only=true]{\fp_sign:n { #1 } }
}
%%%%%%%%%%%%%%%%%%%%%%%%%%%%%%%%%%%%%%%%%%%%%%%%%%%%%%%%%%%%%%%%%%%%%%%%%%%%%%%%
\NewExpandableDocumentCommand{\printFrac}{O{}mm}{%
  %#1: Optional string to trigger display style
  %#2: numerator
  %#3: denominator
  %# Pretty prints the fraction #2/#3
  {
    \ensuremath{
      \str_if_eq:nnT{#1}{d}{\displaystyle}
      \xintTeXsignedFrac{%
        \xintPIrr{\fpeval{#2}/\fpeval{#3}}
      }
    }
  }
}
%%%%%%%%%%%%%%%%%%%%%%%%%%%%%%%%%%%%%%%%%%%%%%%%%%%%%%%%%%%%%%%%%%%%%%%%%%%%%%%%
% Error message for unsupported relations
\msg_new:nnn {exam/signed-ineq}{unknown-relation}
  { The~relation~'#1'~is~not~supported.~Use~<,>,\string\leq,\string\geq. }

% Internal: throw the error
\cs_new_protected:Npn \exam_signed_ineq_error:n #1
  { \msg_error:nnn {exam/signed-ineq}{unknown-relation}{#1} }


% Internal: swap the relation (flip inequality)
\cs_new:Npn \exam_signed_ineq_swap:n #1
  {
    \str_case:nnF {#1}{
      {\leq}{\geq}
      {\geq}{\leq}
      {<}{>}
      {>}{<}
    }{\exam_signed_ineq_error:n {#1}}
  }

% Core: choose swapped or original relation based on the sign of the number
% Uses floating-point comparison so it works for integers and reals.
\cs_new_protected:Npn \exam_signed_ineq:nn #1#2
  {
    \fp_compare:nTF { (#1) < 0 }
      { \exam_signed_ineq_swap:n {#2} }
      { #2 }
  }

% Swap #1 (\geq,\leq,<,>) based on sign of #1
% \signedIneq{<number>}{<relation>}
\NewDocumentCommand{\signedIneq}{ m m }
  { \ensuremath{\exam_signed_ineq:nn {\fp_eval:n{#1}}{#2}} }
%%%%%%%%%%%%%%%%%%%%%%%%%%%%%%%%%%%%%%%%%%%%%%%%%%%%%%%%%%%%%%%%%%%%%%%%%%%%%%%%

% Internal state (unchanged)
\bool_new:N \l_exam_first_bool
\int_new:N  \l_exam_i_int
\int_new:N  \l_exam_start_int
\int_new:N  \l_exam_stop_int
\int_new:N  \l_exam_step_int

% 1) Numeric range: \myRange{start}{stop}{step}
\NewDocumentCommand{\myRange}{mmm}
 { \exam_myrange_numbers_nomathrel:nnn {#1}{#2}{#3} }

\cs_new_protected:Npn \exam_myrange_numbers_nomathrel:nnn #1#2#3
 {
   \int_set:Nn \l_exam_start_int {#1}
   \int_set:Nn \l_exam_stop_int  {#2}
   \int_set:Nn \l_exam_step_int  {#3}
   \bool_set_true:N \l_exam_first_bool

   % If step = 0, do nothing
   \int_if_zero:nTF { \l_exam_step_int } { }
     {
       % Let the stepper handle direction and stopping
       \int_step_inline:nnnn { \l_exam_start_int } { \l_exam_step_int } { \l_exam_stop_int }
         {
           \bool_if:NF \l_exam_first_bool { , }
           \int_to_arabic:n { ##1 }
           \bool_set_false:N \l_exam_first_bool
         }
     }
 }

% 2) Capital letters: \myLetters[<step>]{<StartLetter>}{<StopLetter>}
\NewDocumentCommand{\myLetters}{O{1}mm}
 { \exam_myrange_letters_nomathrel:nnn {#1}{#2}{#3} }

\cs_new_protected:Npn \exam_myrange_letters_nomathrel:nnn #1#2#3
 {
   \int_set:Nn \l_exam_step_int  {#1}
   \int_set:Nn \l_exam_start_int {`#2}  % ASCII code
   \int_set:Nn \l_exam_stop_int  {`#3}
   \bool_set_true:N \l_exam_first_bool

   \int_if_zero:nTF { \l_exam_step_int } { }
     {
       \int_step_inline:nnnn { \l_exam_start_int } { \l_exam_step_int } { \l_exam_stop_int }
         {
           \bool_if:NF \l_exam_first_bool { , }
           \char_generate:nn { ##1 } {12}
           \bool_set_false:N \l_exam_first_bool
         }
     }
 }

\NewDocumentCommand{\myLettersNum}{O{1}mm}
 {
   \exam_myrange_letters_num:nnn {#1}{#2}{#3}
 }

\cs_new_protected:Npn \exam_myrange_letters_num:nnn #1#2#3
 {
   \int_set:Nn \l_exam_step_int  {#1}
   % Convert alphabetical index to ASCII code: 1→65, 2→66, ...
   \int_set:Nn \l_exam_start_int { 64 + #2 }
   \int_set:Nn \l_exam_stop_int  { 64 + #3 }
   \bool_set_true:N \l_exam_first_bool

   % Do nothing if step = 0
   \int_if_zero:nTF { \l_exam_step_int } { }
     {
       \int_step_inline:nnnn { \l_exam_start_int } { \l_exam_step_int } { \l_exam_stop_int }
         {
           \bool_if:NF \l_exam_first_bool { , }
           \char_generate:nn { ##1 } {12}
           \bool_set_false:N \l_exam_first_bool
         }
     }
 }
\ExplSyntaxOn%
%%%%%%%%%%%%%%%%%%%%%%%%%%%%%%%%%%%%%%%%%%%%%%%%%%%%%%%%%%%%%%%%%%%%%%%%%%%%%%%%
% Set random seed based on year, semester, and version
% e.g. VD in spring 2026 corresponds to 1326
% year:           26 : \int_mod:nn { \year }{ 100 }
% exam version   300 : 100 * ( \verNum - 1 )
% semester      1000 : \int_eval:n { \month * 100 + \day

\int_new:N \l__exam_year_two_int
\int_new:N \l__exam_md_int
\int_new:N \l__exam_sem_offset_int
\int_new:N \l__exam_seed_int

% two-digit year
\int_set:Nn \l__exam_year_two_int { \int_mod:nn { \year } { 100 } }

% month-day as MMDD integer (e.g. May 14 -> 514)
\int_set:Nn \l__exam_md_int { \month * 100 + \day }

% semester-based offset:
% before May  5 -> 1000,
% before Jun 20 -> 2000,
% before Aug 15 -> 3000,
% otherwise 4000
\int_case:nnF { 1 }
{
  { \int_compare_p:n { \l__exam_md_int < 505 } }
    { \int_set:Nn \l__exam_sem_offset_int { 1000 } }
  { \int_compare_p:n { \l__exam_md_int < 620 } }
    { \int_set:Nn \l__exam_sem_offset_int { 2000 } }
  { \int_compare_p:n { \l__exam_md_int < 815 } }
    { \int_set:Nn \l__exam_sem_offset_int { 3000 } }
  }
    { \int_set:Nn \l__exam_sem_offset_int { 4000 } }

% base seed = semester offset + two-digit year
\int_set:Nn \l__exam_seed_int { \l__exam_sem_offset_int + \l__exam_year_two_int }

% if a `\version` macro exists, convert to numeric and offset by 100*(ver-1)
\cs_if_exist:NTF \version
  {
    \int_set:Nn \l__exam_seed_int {
      \int_eval:n {
        \l__exam_seed_int + 100 * ( \exp_args:Ne \int_from_alph:n{\version} - 1 ) } }
  }
  { }

% Set global random seed (query with \sys_rand_seed:)
\sys_gset_rand_seed:n { \int_use:N \l__exam_seed_int }
\ExplSyntaxOff%
\ExplSyntaxOff%